\documentclass[10pt,a4paper]{amsart}
\usepackage[margin=19mm]{geometry}
\usepackage{amsmath,amssymb,mathtools}
\usepackage[scr=rsfs,cal=euler]{mathalfa}
\usepackage{xfrac}
\usepackage[unicode,hidelinks,bookmarks=false]{hyperref}
\newcommand{\ie}{i.e.\ }
\newcommand{\cf}{cf.\ }
\newcommand{\resp}{respectively\ }
\newcommand{\Subsection}{\subsection}
\providecommand{\nobreakdash}{\nobreak\hbox{-}}
\setlength{\emergencystretch}{2em}
\multlinegap0pt
\newtheorem{proposition}{Proposition}
\title{Modal analysis of a domain decomposition method for Maxwell's equations in a waveguide}
\author{Victorita Dolean, Antoine Tonnoir, Pierre-Henri Tournier}
\date{Source: arXiv:2502.15548v3 (CC BY 4.0)}
\begin{document}
\maketitle

\noindent\emph{Source and licence.} Modal analysis of a domain decomposition method for Maxwell's equations in a waveguide. Version arXiv:2502.15548v3 (CC BY 4.0), distributed by arXiv under the Creative Commons Attribution 4.0 International licence, http://creativecommons.org/licenses/by/4.0/, as recorded in the arXiv licence metadata for this exact version and shown on the abstract page https://arxiv.org/abs/2502.15548v3. Full licence notice: \texttt{licenses/arxiv-2502.15548-CC-BY-4.0.txt}.\par\smallskip

\noindent\textit{Editorial scope.} Counted unit: the article's proposition prop:BC-TE-TM (source0.tex lines 493--512). The article's complete original proof of it is delivered with it (source0.tex lines 514--608, one \texttt{proof} environment of 95 lines). No mathematics is written, repaired or re-derived: the statement and the proof are the article's own lines, in the article's own notation, and no retained line is substituted or re-flowed.  No further source-local context is needed: every cross-reference of every retained line resolves inside the two delivered spans.  Nothing was omitted from any retained span, so the package carries no omission record.  No retained line contains a citation, so the wrapper carries no bibliography and no bibliography entry.  No retained line contains a figure environment, an image-inclusion command or drawing code, so no image is delivered and the package declares no asset dependency.  Editorial formatting only, in addition: an AMS \texttt{amsart} preamble that restates the article's own declarations and macros used by the retained lines, namely the environments \texttt{proposition} on separate counters, together with the standard packages those lines need.  The retained environments print in this document as Proposition 1 in order of appearance, whereas the article prints the same items with the numbering of its own full document.  The complete native source of the article as distributed in the arXiv e-print for this exact version is bundled unchanged as \texttt{sources/173656/source0.tex}.
\begin{proposition}[Diagonalization of the tangential curl operator]\label{prop:BC-TE-TM}
Let $\mathbf n=(1,0,0)^t$ be the normal to a cross-section of the waveguide. Then the operator
\[
\mathbf E \;\longmapsto\; ((\nabla\times \mathbf E)\times \mathbf n)\times \mathbf n
\]
acts diagonally on TE, TM, and TEM modes. More precisely, for every modal index $i$,
\[
((\nabla\times \mathbf E_i^{{\rm T},\pm})\times \mathbf n)\times \mathbf n
=
\mu_i^{\rm T,\pm}\,(\mathbf E_i^{{\rm T},\pm}\times \mathbf n),
\]
where
\[
\mu_i^{\rm TE,\pm}=\pm \imath \beta_i^{\rm TE},
\qquad
\mu_i^{\rm TM,\pm}=\pm \imath \frac{k^2}{\beta_i^{\rm TM}},
\qquad
\mu_i^{\rm TEM,\pm}=\pm \imath k.
\]
\end{proposition}

\begin{proof}
For any field $\mathbf E=(E_x,E_y,E_z)^t$, one has
\begin{equation}\label{eq:curl-cross-n}
(\nabla\times \mathbf E)\times \mathbf n
=
\begin{bmatrix}
0\\
\partial_x E_y-\partial_y E_x\\
\partial_x E_z-\partial_z E_x
\end{bmatrix}.
\end{equation}

\smallskip
\noindent\textbf{TE and TEM modes.}
For TE modes, $E_x^{\rm TE,\pm}=0$. Since
\[
\mathbf E^{\rm TE,\pm}_i
=
e^{\pm \imath \beta_i^{\rm TE}x}\widehat{\mathbf E}_i^{\rm TE,\pm},
\]
formula \eqref{eq:curl-cross-n} gives
\[
(\nabla\times \mathbf E_i^{\rm TE,\pm})\times \mathbf n
=
\pm \imath \beta_i^{\rm TE}\,\mathbf E_i^{\rm TE,\pm}.
\]
Taking once more the cross product with $\mathbf n$ yields
\[
((\nabla\times \mathbf E_i^{\rm TE,\pm})\times \mathbf n)\times \mathbf n
=
\pm \imath \beta_i^{\rm TE}\,(\mathbf E_i^{\rm TE,\pm}\times \mathbf n).
\]
The same argument applies to TEM modes, since their longitudinal electric component also vanishes and $\beta_i^{\rm TEM}=k$.

\smallskip
\noindent\textbf{TM modes.}
For TM modes, one has $H_x^{\rm TM,\pm}=0$ and
\[
\mathbf H_i^{\rm TM,\pm}
=
e^{\pm \imath \beta_i^{\rm TM}x}\widehat{\mathbf H}_i^{\rm TM,\pm}.
\]
Applying \eqref{eq:curl-cross-n} to $\mathbf H_i^{\rm TM,\pm}$ gives
\[
(\nabla\times \mathbf H_i^{\rm TM,\pm})\times \mathbf n
=
\pm \imath \beta_i^{\rm TM}\,\mathbf H_i^{\rm TM,\pm}.
\]
Using the first-order Maxwell equation
\[
\nabla\times \mathbf H_i^{\rm TM,\pm}
=
-\imath \omega\varepsilon\,\mathbf E_i^{\rm TM,\pm},
\]
we obtain
\[
-\imath \omega\varepsilon\,(\mathbf E_i^{\rm TM,\pm}\times \mathbf n)
=
\pm \imath \beta_i^{\rm TM}\,\mathbf H_i^{\rm TM,\pm}.
\]
Now, from
\[
\nabla\times \mathbf E_i^{\rm TM,\pm}
=
\imath \omega\mu\,\mathbf H_i^{\rm TM,\pm},
\]
it follows that
\[
\nabla\times \mathbf E_i^{\rm TM,\pm}
=
\mp \imath \frac{k^2}{\beta_i^{\rm TM}}\,
(\mathbf E_i^{\rm TM,\pm}\times \mathbf n).
\]
Taking the tangential trace gives
\[
((\nabla\times \mathbf E_i^{\rm TM,\pm})\times \mathbf n)\times \mathbf n
=
\mp \imath \frac{k^2}{\beta_i^{\rm TM}}
\bigl((\mathbf E_i^{\rm TM,\pm}\times \mathbf n)\times \mathbf n\bigr)\times \mathbf n.
\]
Finally, since
\[
\bigl((\mathbf u\times \mathbf n)\times \mathbf n\bigr)\times \mathbf n
=
-\mathbf u\times \mathbf n
\]
for any vector field $\mathbf u$, we conclude that
\[
((\nabla\times \mathbf E_i^{\rm TM,\pm})\times \mathbf n)\times \mathbf n
=
\pm \imath \frac{k^2}{\beta_i^{\rm TM}}\,
(\mathbf E_i^{\rm TM,\pm}\times \mathbf n).
\]
This proves the result.
\end{proof}

\end{document}
